\documentclass[11pt]{article}

\usepackage[T1]{fontenc}
\usepackage{lmodern}
\usepackage{microtype}
\usepackage[a4paper,margin=25mm,headheight=14pt]{geometry}
\usepackage{amsmath,amssymb,amsthm,mathtools}
\usepackage{booktabs,longtable,array,tabularx}
\usepackage{enumitem}
\usepackage{xcolor}
\usepackage{pid-rs-report-tables}
\usepackage{hyperref}
\usepackage{xurl}
\usepackage{fancyhdr}
\usepackage{listings}

\setlength{\emergencystretch}{3em}
\newcolumntype{L}[1]{>{\raggedright\arraybackslash}p{#1}}

\ifdefined\pdfinfoomitdate
  \pdfinfoomitdate=1
\fi
\ifdefined\pdftrailerid
  \pdftrailerid{}
\fi
\ifdefined\pdfsuppressptexinfo
  \pdfsuppressptexinfo=15
\fi

\hypersetup{%
  colorlinks=true,
  linkcolor=blue!55!black,
  citecolor=blue!55!black,
  urlcolor=blue!55!black,
  pdftitle={Exact Two-Source Categorical SxPID Count-to-Atom Bridge},
  pdfsubject={Bounded formalization of the 24 informative, misinformative, and signed-net cumulative and atom coordinates},
  pdfkeywords={partial information decomposition, shared exclusions, Mobius inversion, exact counts, Lean, formal verification},
  pdfauthor={pid-rs project analysis}
}

\lstset{%
  basicstyle=\ttfamily\small,
  breaklines=true,
  columns=fullflexible,
  frame=single,
  framerule=0.3pt,
  xleftmargin=0.5em,
  xrightmargin=0.5em,
  aboveskip=0.7em,
  belowskip=0.7em,
  showstringspaces=false
}

\PidConfigureSections
\PidSetRunningHeads{pid-rs two-source count-to-atom bridge}{bounded formal audit}
\PidConfigureAbstract
\PidConfigureLinksAndListings

\newtheorem{theorem}{Theorem}[section]
\newtheorem{corollary}[theorem]{Corollary}
\newtheorem{proposition}[theorem]{Proposition}
\newtheorem{lemma}[theorem]{Lemma}
\theoremstyle{definition}
\newtheorem{definition}[theorem]{Definition}
\newtheorem{remark}[theorem]{Remark}

\newcommand{\sx}{\mathrm{sx}}
\newcommand{\supp}{\operatorname{supp}_{+}}
\newcommand{\Avg}{\operatorname{Avg}}
\newcommand{\Count}{\operatorname{Count}}
\newcommand{\Mob}{\operatorname{Mob}}
\newcommand{\Zeta}{\operatorname{Zeta}}

\title{\PidReportTitle
  {Exact Two-Source Categorical SxPID\\Count-to-Atom Bridge}
  {A bounded 24-coordinate Lean formalization of components, M\"obius atoms,\\
   exact count products, and sign reduction}}
\author{pid-rs project analysis}
\date{10 August 2026}

\begin{document}
\PidMakeTitle

\begin{abstract}
This report documents the bounded Lean formalization for the complete two-source finite
categorical shared-exclusions surface: four cumulative nodes and four M\"obius atoms, each carrying
informative, misinformative, and signed-net components.  The checked module fixes an explicit
24-coordinate order, proves the concrete two-source M\"obius and zeta transforms are inverse,
proves source-labelled coordinate-swap equivariance, and proves that finite empirical averaging commutes with the
transform.  Starting from an arbitrary supplied natural-valued count function with positive total,
it relates every supported local component to an exact rational logarithm argument.  Every averaged
coordinate is then represented as \(N^{-1}\log Q\) for a positive exact rational product \(Q\), so
positive, negative, and zero values reduce to comparison of \(Q\) with one.

The categorical event functional, component split, and M\"obius atom construction are defined in
the Makkeh--Gutknecht--Wibral publications.  The Lean encodings, coordinate order, supplied-count
interface, exact product layer, theorem factoring, and assurance packet are project-defined
formalization work; they define no new PID and carry no scientific-priority claim.  Exact pinned
kernel replay, 80 registered static and baseline-first isolated changes, categorical-only scope
binding, deterministic PDF checks, page inspection, and separate read-only repository-workflow
reviews are recorded.  Component-atom nonnegativity, Rust and binary64 refinement,
standalone-certifier refinement, statistical conclusions, and higher-source lattices are explicit
nonclaims.
\end{abstract}

\tableofcontents

\section{Status and decision boundary}

\noindent\colorbox{PidSoft}{%
  \parbox{\dimexpr\linewidth-2\fboxsep\relax}{%
    \textbf{Bounded accepted status.}  The exact source is accepted only for the supplied-count,
    fixed-two-source categorical formal scope.  It is not a publication-correspondence,
    component-sign, executable-refinement, higher-source, or population theorem.}}

The checked source is \path{TwoSourceMobiusAtomBridge.lean} under
\path{audit/formal/lean/PidFiniteConvergence}.  It imports the completed supplied-count event bridge
and adds the concrete two-source component/atom layer.  The accepted target is exactly
\begin{equation}\label{eq:coordinate-count}
3\times(4+4)=24.
\end{equation}

The three components are informative, misinformative, and signed net.  The four cumulative nodes
are source one, source two, joint sources, and redundancy.  The four atoms are unique one, unique
two, synergy, and redundancy.  Equation~\eqref{eq:coordinate-count} is a finite categorical
surface, not a continuous estimator and not a higher-source lattice.

The acceptance evidence comprises the following bounded layers:
\begin{enumerate}[leftmargin=*]
  \item pinned Lean toolchain, dependency, source, declaration, and axiom-basis replay;
  \item meaningful negative semantic variations of order, events, coefficients, components,
        support, weights, products, comparisons, and scope prose;
  \item method-catalog, assurance-registry, formal-map, limitations, and changelog integration;
  \item deterministic PDF generation, extraction, font and geometry checks, and page inspection;
        and
  \item separate read-only mathematical, provenance, formal, implementation-boundary, and
        presentation reviews within the repository workflow.
\end{enumerate}

\section{Scientific provenance and project contribution}

\subsection{Paper-defined object}

Makkeh, Gutknecht, and Wibral define the finite categorical pointwise shared-exclusions
construction, its informative and misinformative decomposition, the signed-net quantity, and its
partial-information atoms through M\"obius inversion~\cite{mgw2021}.  Gutknecht, Wibral, and
Makkeh develop the associated part-whole and formal-logic organization~\cite{gwm2021}.

The following scientific content is therefore paper-defined:
\begin{itemize}[leftmargin=*]
  \item keyed source-event disjunctions and keyed target restrictions;
  \item pointwise informative and misinformative cumulative components;
  \item signed net as informative minus misinformative;
  \item the passage from cumulatives to atoms by M\"obius inversion; and
  \item empirical averaging of pointwise quantities.
\end{itemize}

\subsection{Project-defined formalization}

The repository contributes a formal and assurance realization of that fixed object:
\begin{itemize}[leftmargin=*]
  \item heterogeneous finite source types indexed by \(\mathrm{Fin}\,2\) and a finite target type;
  \item explicit Lean constructors for components, nodes, atoms, and cumulative-or-atom
        coordinates;
  \item fixed node, atom, component, and complete 24-coordinate orders;
  \item concrete integer M\"obius and zeta coefficient tables;
  \item exact supplied-count component arguments on positive support;
  \item exact rational and real product representatives for every coordinate;
  \item generic additive-group inverse and reconstruction theorems;
  \item theorem-level links between pointwise averaging, cumulative inversion, count formulas,
        products, and signs; and
  \item a bounded replay and review plan.
\end{itemize}

These are project-defined encodings, proofs, and engineering boundaries.  They do not define a new
PID measure.  The correspondence from publication equations to the selected Lean definitions is a
human-reviewed transcription boundary, not a theorem that Lean derives from publication text.

\section{Inherited supplied-count event model}

\subsection{Complete keys and positive support}

Let the two source alphabets \(\mathcal S_1,\mathcal S_2\) and target alphabet
\(\mathcal T\) be finite.  A complete key is
\[
z=(s_1,s_2,t)\in\mathcal S_1\times\mathcal S_2\times\mathcal T.
\]
Supply an arbitrary natural count function \(c\) and define
\begin{equation}\label{eq:total}
N=\sum_z c(z)>0,
\qquad
Z_+=\{z:c(z)>0\}.
\end{equation}
Zero-count keys are permitted.  The exact empirical law is
\begin{equation}\label{eq:empirical-law}
\widehat p_c(z)=\frac{c(z)}{N},
\end{equation}
and all local logarithms and empirical sums below range over \(Z_+\).

\subsection{Keyed events}

For a cumulative node \(\alpha\) and supported anchor \(z\), let
\begin{align}
E_\alpha(z)&=\text{the keyed source event},\\
T(z)&=\text{the keyed target branch},\\
E_{\alpha,t}(z)&=E_\alpha(z)\cap T(z).
\end{align}
Write their exact natural counts as
\begin{equation}\label{eq:event-counts}
C_\alpha(z)=\Count(E_\alpha(z)),
\quad
C_t(z)=\Count(T(z)),
\quad
C_{\alpha,t}(z)=\Count(E_{\alpha,t}(z)).
\end{equation}
The completed count/event dependency proves that a supported anchor belongs to each defining
event.  Consequently
\begin{equation}\label{eq:count-positive}
C_\alpha(z)>0,
\qquad C_t(z)>0,
\qquad C_{\alpha,t}(z)>0
\qquad(z\in Z_+).
\end{equation}
Positivity in equation~\eqref{eq:count-positive} is derived from anchor membership; it is not an
extra statistical assumption.

\section{Fixed cumulative nodes and components}

\subsection{Four cumulative nodes}

\begin{table}[ht]
\centering
\caption{Exact two-source cumulative-node order.}\label{tab:nodes}
\begin{tabularx}{\linewidth}{@{}rL{0.22\linewidth}XL{0.12\linewidth}@{}}
\toprule
\PidTableHeaderRow
Position & Lean node & Keyed source-event meaning & Symbol \\
\midrule
1 & \texttt{sourceOne} & Source one matches its anchor value. & \(C_1\) \\
2 & \texttt{sourceTwo} & Source two matches its anchor value. & \(C_2\) \\
3 & \texttt{jointSources} & Both sources match their anchor values. & \(C_{12}\) \\
4 & \texttt{redundancy} & Source one or source two matches its anchor value. & \(C_R\) \\
\bottomrule
\end{tabularx}
\end{table}

The redundancy event in Table~\ref{tab:nodes} is a union; the joint-source event is the
intersection of the two singleton source branches.  Cumulative source-one and source-two values are
not unique atoms, and the joint-source cumulative is not the synergy atom.

\subsection{Three local component arguments}

For \(u\in\{+,-,n\}\), define the exact rational local arguments
\begin{align}
q^+_\alpha(z)&=\frac{N}{C_\alpha(z)},\label{eq:qplus}\\
q^-_\alpha(z)&=\frac{C_t(z)}{C_{\alpha,t}(z)},\label{eq:qminus}\\
q^n_\alpha(z)&=\frac{N C_{\alpha,t}(z)}{C_\alpha(z)C_t(z)}.
\label{eq:qnet}
\end{align}
Equation~\eqref{eq:count-positive} makes every argument strictly positive.  Direct rational
algebra gives
\begin{equation}\label{eq:q-ratio}
q^n_\alpha(z)=\frac{q^+_\alpha(z)}{q^-_\alpha(z)}.
\end{equation}

The local component values in nats are
\begin{equation}\label{eq:local-components}
i^u_\alpha(z)=\log q^u_\alpha(z).
\end{equation}
The logarithm quotient rule and equation~\eqref{eq:q-ratio} give
\begin{equation}\label{eq:local-net}
i^n_\alpha(z)=i^+_\alpha(z)-i^-_\alpha(z).
\end{equation}

The checked module proves the informative and misinformative empirical probability expressions equal
the logarithms in equations~\eqref{eq:qplus} and~\eqref{eq:qminus}; the inherited module already
supplies the net case.  A single component selector then states the result uniformly.

\subsection{Averaged cumulatives}

For each component and node, define
\begin{equation}\label{eq:averaged-cumulative}
C^u_\alpha
=\sum_{z\in Z_+}\frac{c(z)}{N}i^u_\alpha(z).
\end{equation}
Finite-sum linearity yields
\begin{equation}\label{eq:averaged-net}
C^n_\alpha=C^+_\alpha-C^-_\alpha.
\end{equation}
Equation~\eqref{eq:averaged-cumulative} uses exact count weights.  It inserts no pseudo-count,
smoothing rule, uniform-on-support weight, clipping, or absolute value.

\section{Concrete two-source M\"obius algebra}

\subsection{Coefficient tables}

Use cumulative column order \([C_1,C_2,C_{12},C_R]\) and atom row order
\([U_1,U_2,S,R]\).  The checked module fixes
\begin{equation}\label{eq:mobius-matrix}
M=
\begin{pmatrix}
 1& 0&0&-1\\
 0& 1&0&-1\\
-1&-1&1& 1\\
 0& 0&0& 1
\end{pmatrix}.
\end{equation}
The concrete M\"obius transform is therefore
\begin{align}
\Pi^u_R&=C^u_R,\label{eq:atom-r}\\
\Pi^u_{U_1}&=C^u_1-C^u_R,\label{eq:atom-u1}\\
\Pi^u_{U_2}&=C^u_2-C^u_R,\label{eq:atom-u2}\\
\Pi^u_S&=C^u_{12}-C^u_1-C^u_2+C^u_R.\label{eq:atom-s}
\end{align}

The inverse zeta coefficient matrix is
\begin{equation}\label{eq:zeta-matrix}
Z=
\begin{pmatrix}
1&0&0&1\\
0&1&0&1\\
1&1&1&1\\
0&0&0&1
\end{pmatrix},
\end{equation}
with cumulative row order and atom column order as above.

\begin{proposition}[Inverse algebra]\label{prop:inverse}
For any additive commutative group, the concrete transforms satisfy
\[
\Zeta(\Mob(C))=C,
\qquad
\Mob(\Zeta(\Pi))=\Pi.
\]
\end{proposition}

The source proves Proposition~\ref{prop:inverse} by cases on the four positions and commutative
group algebra.  The theorem is independent of real numbers, logarithms, and count semantics.

\begin{corollary}[Reconstruction identities]\label{cor:reconstruct}
For each component,
\begin{align}
C^u_R&=\Pi^u_R,\\
C^u_1&=\Pi^u_{U_1}+\Pi^u_R,\\
C^u_2&=\Pi^u_{U_2}+\Pi^u_R,\\
C^u_{12}&=\Pi^u_{U_1}+\Pi^u_{U_2}+\Pi^u_S+\Pi^u_R.
\end{align}
\end{corollary}

\subsection{Integer-row representation}

The checked module separately states that each concrete transform equals the finite sum of its integer
coefficient row.  This matters because the coefficient table, rather than only a simplified
formula, is part of the intended formal contract.  It also proves that every nonredundancy
M\"obius row sums to zero while the redundancy row sums to one.

\subsection{Source-labelled coordinate exchange}

Define coordinate exchange on source-labelled cumulatives and atoms by
\begin{equation}\label{eq:swap}
C_1\leftrightarrow C_2,
\qquad U_1\leftrightarrow U_2,
\end{equation}
while joint sources, redundancy, synergy, and redundancy atom remain fixed.  The checked module proves
both maps are involutions and
\begin{equation}\label{eq:equivariant}
\Mob(C\circ\operatorname{swap})(\gamma)
=\Mob(C)(\operatorname{swap}(\gamma)).
\end{equation}
This is formula equivariance for an arbitrary four-entry cumulative function.  It does not
transport heterogeneous source types, keys, counts, laws, or inherited events, and it does not
relate the order to Rust fields or a serialized record.

\section{Averaging and component linearity}

\subsection{Pointwise atoms}

Let \(m_{\gamma\alpha}\) be the integer coefficient in row \(\gamma\) of
equation~\eqref{eq:mobius-matrix}.  Define a pointwise atom by
\begin{equation}\label{eq:pointwise-atom}
\pi^u_\gamma(z)=\sum_\alpha m_{\gamma\alpha}i^u_\alpha(z).
\end{equation}
Its positive-support empirical average is
\begin{equation}\label{eq:average-pointwise}
\overline\pi^u_\gamma
=\sum_{z\in Z_+}\frac{c(z)}{N}\pi^u_\gamma(z).
\end{equation}

\begin{proposition}[Averaging/M\"obius commutation]\label{prop:commute}
For every component and atom,
\begin{equation}\label{eq:commute}
\overline\pi^u_\gamma
=\sum_\alpha m_{\gamma\alpha}C^u_\alpha
=\Pi^u_\gamma.
\end{equation}
\end{proposition}

The proof expands the four atom cases and distributes multiplication and finite sums over addition
and subtraction.  It establishes equality of ``invert each supported point then average'' and
``average each cumulative then invert'' for the selected finite empirical formalization.

\subsection{Net component}

The M\"obius transform is additive and commutes with subtraction.  Combining this fact with
equation~\eqref{eq:averaged-net} yields
\begin{equation}\label{eq:atom-net}
\Pi^n_\gamma=\Pi^+_\gamma-\Pi^-_\gamma.
\end{equation}
Equation~\eqref{eq:atom-net} is an exact signed identity.  No intermediate term or final atom is
silently clamped.

\section{The exact 24-coordinate surface}

\subsection{Types and order}

The checked module defines three finite types:
\begin{itemize}[leftmargin=*]
  \item component: informative, misinformative, net;
  \item atom: unique one, unique two, synergy, redundancy; and
  \item coordinate: a component-indexed cumulative node or a component-indexed atom.
\end{itemize}

It fixes the explicit order in Table~\ref{tab:coordinates}.

\begin{longtable}{@{}L{0.30\linewidth}L{0.62\linewidth}@{}}
\caption{Project-defined order over the 24 paper-defined quantities.}\label{tab:coordinates}\\
\toprule
\PidTableHeaderRow
Block & Ordered coordinates \\
\midrule
\endfirsthead
\toprule
\PidTableHeaderRow
Block & Ordered coordinates \\
\midrule
\endhead
Informative cumulatives & \(C^+_1,C^+_2,C^+_{12},C^+_R\) \\
Misinformative cumulatives & \(C^-_1,C^-_2,C^-_{12},C^-_R\) \\
Signed-net cumulatives & \(C^n_1,C^n_2,C^n_{12},C^n_R\) \\
Informative atoms & \(\Pi^+_{U_1},\Pi^+_{U_2},\Pi^+_S,\Pi^+_R\) \\
Misinformative atoms & \(\Pi^-_{U_1},\Pi^-_{U_2},\Pi^-_S,\Pi^-_R\) \\
Signed-net atoms & \(\Pi^n_{U_1},\Pi^n_{U_2},\Pi^n_S,\Pi^n_R\) \\
\bottomrule
\end{longtable}

The source states that the coordinate type has cardinality 24, the ordered list has length 24 with
no duplicate, and the list's finite-set projection equals the complete coordinate universe.  This
fixes a mathematical order.  It does not prove equality to a Rust array, report field sequence,
or external encoding.

\subsection{Uniform count-expression theorem}

Define a count expression for a cumulative by inserting equations~\eqref{eq:qplus}--\eqref{eq:qnet}
into equation~\eqref{eq:averaged-cumulative}.  Define the atom count expression by applying the
same concrete M\"obius transform.  The checked module's central theorem is quantified over the
coordinate type:
\begin{lstlisting}
theorem all_24_averaged_coordinates_empirical_eq_count_expression
    (count : CategoricalKey (Fin 2) sourceValue targetValue -> Nat)
    (coordinate : SxPid2Coordinate)
    (h_total : 0 < totalCount count) :
    averagedSxPid2Coordinate (empiricalLaw count) coordinate =
      sxPid2CountCoordinateExpression count coordinate
\end{lstlisting}

The cumulative branch uses the supported local empirical-log theorem.  The atom branch transports
the four cumulative equalities through the M\"obius transform.

\section{Exact rational products}

\subsection{Cumulative products}

For each component and cumulative node, define
\begin{equation}\label{eq:cumulative-product}
P^u_\alpha
=\prod_{z\in Z_+}\left(q^u_\alpha(z)\right)^{c(z)}.
\end{equation}
Every factor is a positive rational, and each exponent is a natural count.  The checked module defines
both an exact rational value and its real counterpart, and proves that casting the exact rational
product gives the real product.

\begin{proposition}[Cumulative normalization]\label{prop:cumulative-product}
For every component and cumulative node,
\begin{equation}\label{eq:cumulative-scaled-log}
C^u_\alpha=\frac1N\log P^u_\alpha.
\end{equation}
\end{proposition}

The proof applies the logarithm of a finite product and the logarithm of a natural power, then
factors the positive scalar \(1/N\) out of the weighted sum.

\subsection{Atom products}

Write \(A^u_\gamma\) for an atom product and \(P^u_\alpha\) for a cumulative product.  The
multiplicative form of equations~\eqref{eq:atom-r}--\eqref{eq:atom-s} is
\begin{align}
A^u_{U_1}&=\frac{P^u_1}{P^u_R},\label{eq:product-u1}\\
A^u_{U_2}&=\frac{P^u_2}{P^u_R},\label{eq:product-u2}\\
A^u_S&=\frac{P^u_{12}P^u_R}{P^u_1P^u_2},\label{eq:product-s}\\
A^u_R&=P^u_R.\label{eq:product-r}
\end{align}
All cumulative products are positive, so every quotient is defined and positive.

\begin{proposition}[Atom normalization]\label{prop:atom-product}
For every component and atom,
\begin{equation}\label{eq:atom-scaled-log}
\Pi^u_\gamma=\frac1N\log A^u_\gamma.
\end{equation}
\end{proposition}

The proof uses equations~\eqref{eq:cumulative-scaled-log} and
\eqref{eq:product-u1}--\eqref{eq:product-r}, together with logarithm identities justified by
strict positivity.

\subsection{Uniform coordinate normalization}

The cumulative and atom results combine into one statement:
\begin{equation}\label{eq:all-coordinate-log}
V(c,\xi)=\frac1N\log Q(c,\xi)
\qquad\text{for every one of the 24 coordinates }\xi.
\end{equation}
The checked module also proves
\begin{equation}\label{eq:rational-cast}
Q_{\mathbb R}(c,\xi)
=\bigl(Q_{\mathbb Q}(c,\xi):\mathbb R\bigr),
\end{equation}
where \(Q_{\mathbb Q}\) is an exact rational value.

\section{Exact sign and zero reduction}

Equation~\eqref{eq:total} gives \(1/N>0\), and product positivity gives
\(Q(c,\xi)>0\).  The strict
monotonicity of the natural logarithm on positive reals and \(\log 1=0\) yield the following.

\begin{theorem}[All-coordinate sign reduction]\label{thm:sign}
For every supplied count function with positive total and every coordinate \(\xi\),
\begin{align}
V(c,\xi)>0&\iff Q(c,\xi)>1,\label{eq:positive-iff}\\
V(c,\xi)<0&\iff Q(c,\xi)<1,\label{eq:negative-iff}\\
V(c,\xi)=0&\iff Q(c,\xi)=1.\label{eq:zero-iff}
\end{align}
\end{theorem}

Theorem~\ref{thm:sign} removes transcendental approximation from the mathematical classification
once the exact rational product is available.  It does not prove that any current parser,
standalone certifier, Rust implementation, or report produces the same rational product or performs
the comparison with a declared resource policy.

\subsection{Component-nonnegativity is not proved}

Equations~\eqref{eq:positive-iff}--\eqref{eq:zero-iff} are conditional equivalences.  They tell us
how a coordinate's sign follows from the corresponding product.  They do not establish where the
product for every informative or misinformative atom must lie.  In particular, the checked module does
not prove
\[
\Pi^+_\gamma\ge0
\qquad\text{or}\qquad
\Pi^-_\gamma\ge0
\quad\text{for all }\gamma.
\]
Nor does it produce a counterexample to either statement.  The publication's component-sign
results must not be reported as Lean-proved by this packet unless a separate formal theorem is
added and accepted.

\section{Theorem crosswalk}

\begin{longtable}{@{}L{0.29\linewidth}L{0.63\linewidth}@{}}
\caption{Main bound declaration groups.}\label{tab:theorems}\\
\toprule
\PidTableHeaderRow
Layer & Bound declarations \\
\midrule
\endfirsthead
\toprule
\PidTableHeaderRow
Layer & Bound declarations \\
\midrule
\endhead
Types and order & \texttt{SxPid2Component}, \texttt{SxPid2Atom},
\texttt{SxPid2Coordinate}, the four order definitions, and their length, no-duplicate, and
completeness/cardinality theorems. \\
Transforms & \texttt{sxPid2MobiusCoefficient}, \texttt{sxPid2ZetaCoefficient},
\texttt{sxPid2MobiusTransform}, \texttt{sxPid2ZetaTransform}, and their integer-row forms. \\
Inverse algebra & \texttt{sx\_pid2\_zeta\_after\_mobius},
\texttt{sx\_pid2\_mobius\_after\_zeta}, and cumulative reconstruction theorems. \\
Source exchange & Node and atom swap definitions, involution theorems, and
\texttt{sx\_pid2\_mobius\_coordinate\_swap\_equivariant}. \\
Components & Local, averaged-cumulative, local-atom, averaged-pointwise-atom, and averaged-atom
component selectors, plus the net-subtraction theorems. \\
Averaging & \path{averaged_pointwise_atom_eq_mobius_of_averaged_cumulatives}. \\
Local counts & Informative, misinformative, and selected component arguments; positivity; net
ratio; and local empirical-log equalities. \\
All-coordinate counts & Cumulative, atom, and coordinate count expressions, culminating in
\path{all_24_averaged_coordinates_empirical_eq_count_expression}. \\
Products & Cumulative, atom, and coordinate real/rational products, positivity, and exact-rational
cast equalities. \\
Scaled logarithms & Cumulative and atom normalization, then
\texttt{all\_24\_averaged\_coordinates\_eq\_scaled\_log\_product}. \\
Signs & The all-coordinate positive, negative, and zero equivalences. \\
\bottomrule
\end{longtable}

This table is a human crosswalk, not the enforcement route.  The checker binds exact source bytes,
339 declarations, 246 named source theorem axiom bases, and both semantic contracts.

\section{Twenty-four-lens bounded review}

\begin{longtable}{@{}rL{0.19\linewidth}L{0.66\linewidth}@{}}
\caption{Structured review of the accepted boundary.}\label{tab:lenses}\\
\toprule
\PidTableHeaderRow
No. & Lens & Bounded finding \\
\midrule
\endfirsthead
\toprule
\PidTableHeaderRow
No. & Lens & Bounded finding \\
\midrule
\endhead
1 & Measure identity & The definitions remain finite categorical shared exclusions.  No
continuous estimator, \(I_{\min}\), or other PID is substituted. \\
2 & Publication provenance & Event, component, and M\"obius mathematics is paper-defined; the
Lean realization and assurance packet are project-defined. \\
3 & Source arity & Every new scientific coordinate is fixed to two sources through
\(\mathrm{Fin},2\). \\
4 & Alphabet generality & The two finite source alphabets may be heterogeneous; the target is
finite. \\
5 & Count quantifier & Counts are arbitrary natural values and need not be positive at every
complete key. \\
6 & Total and support & Positive total is explicit; logarithms and averages range only over
positive-count anchors. \\
7 & Event dependency & The module imports the completed keyed event semantics rather than
silently redefining source or target events. \\
8 & Component relation & Signed net is informative minus misinformative, with no clamping or
absolute value. \\
9 & Node order & Source one, source two, joint, and redundancy have an explicit order and distinct
constructors. \\
10 & Atom order & Unique one, unique two, synergy, and redundancy are explicit and are not
confused with cumulative nodes. \\
11 & Coefficient visibility & Integer M\"obius and zeta tables are definitions rather than
unstated algebra. \\
12 & Two-sided inverse & Both transform compositions are theorem targets over an additive
commutative group. \\
13 & Reconstruction & Marginal and joint cumulative reconstruction equations are explicit. \\
14 & Coordinate symmetry & Swapping source-labelled coordinate positions exchanges only the
source-specific node and atom pairs; no typed source-data transport is claimed. \\
15 & Component linearity & M\"obius inversion commutes with informative-minus-misinformative
subtraction. \\
16 & Averaging order & Pointwise inversion then averaging equals averaging cumulatives then
inverting. \\
17 & Coordinate completeness & Coordinate cardinality, explicit-list length, no-duplicate, and
\(
\operatorname{toFinset}=\operatorname{univ}
\) facts all establish the same 24-coordinate universe. \\
18 & Local exactness & Each supported local component is related to an exact rational count
argument. \\
19 & Empirical weighting & Every local logarithm is weighted by its exact multiplicity divided by
total count. \\
20 & Product construction & Multiplicities become natural exponents; negative M\"obius
coefficients become exact quotients. \\
21 & Exact-real bridge & Real products are tied to casts of exact rational products. \\
22 & Sign logic & Strict positivity, strict negativity, and exact zero use comparison with one and
the positive scale \(1/N\). \\
23 & Refinement boundary & Rust, binary64, parsing, standalone-certifier, Python, and resource
semantics remain outside the theorem. \\
24 & Acceptance evidence & Exact kernel/source binding, 80 registered changes, categorical-only
scope integration, deterministic PDF checks, page inspection, and separate read-only
repository-workflow reviews are recorded. \\
\bottomrule
\end{longtable}

The lenses in Table~\ref{tab:lenses} are one structured review, not 24 independent proofs or 24
independent implementations.

\section{Residual boundaries}

The accepted bounded claim must not expand
to any of the following:
\begin{enumerate}[leftmargin=*]
  \item derivation of Lean definitions from publication text;
  \item row, byte, file, JSON, or other representation to the supplied count function;
  \item Rust sorting, histogram extraction, or event scanning;
  \item Rust \texttt{NODES2}, \texttt{invert2}, pointwise accumulation, atom fields, or order;
  \item binary64, MPFR, Python, compiler, runtime, or hardware equality;
  \item standalone-certifier schemas, parsers, product evaluation, comparisons, or reports;
  \item integer overflow, expression normalization, allocation, cancellation, time, memory, or
        another resource theorem;
  \item informative or misinformative atom nonnegativity;
  \item sampling uncertainty, concentration, consistency, population support, or calibration;
  \item continuous Ehrlich/KSG shared exclusions, fitted quantized wrappers, \(I_{\min}\), or
        Shannon-invariant families;
  \item three-source or general higher-source lattices;
  \item scientific priority, uniqueness, application validity, release readiness, or downstream
        authority.
\end{enumerate}

Executable agreement on bounded fixtures can be useful corroboration after it is explicitly
performed.  It does not replace a missing formal refinement theorem.  Likewise, exact finite-table
algebra does not imply a sampling or population conclusion.

\section{Acceptance evidence}

The pinned checker binds 11 Lean sources, eight imported modules, 339 declarations, and 246 named
source theorem axiom bases.  The original 16-example and atom 11-example semantic contracts compile
with \texttt{lean -t 0}; the atom contract also checks the permitted basis of six private helper
theorems.  Under normal and optimized Python, the self-test rejects 49 static, five count/event,
17 atom-module, and nine atom-contract changes.  The six previously missing minimum classes cover
same-name inverse weakening, coordinate omission and duplication, net addition, atom-product
quotient-to-multiplication, product-comparison reversal, and removal of a logarithm-product
positivity premise.

The method catalog and generated assurance registry route the atom-exclusive evidence only to
the stable categorical SxPID family.  The Markdown/LaTeX/PDF pair is checked for deterministic rendering,
strict LaTeX diagnostics, required boundary text, A4 geometry, and embedded Unicode-mapped fonts.
Every rendered page was inspected for clipping, overlap, table continuation, and legibility.
Separate read-only formal, integration, and document reviews within the repository workflow
checked the mathematics, provenance, executable boundary, checker sensitivity, and presentation.
These are not external review or independent authorship.  The exact current hashes, commands, and
outcomes are recorded in the versioned Lean 4.33 replay receipt.  The dated phase-A receipt remains
the immutable historical Lean 4.32 observation of the earlier 71-route suite; it is not evidence
for the current 80-route replay.  Both receipts and the residual cut are linked from the canonical
Markdown audit.

\section{Conclusion}

The checked module closes a precise mathematical gap between the completed two-source count/event
semantics and a complete finite categorical component/atom surface.  It makes the two-source
M\"obius algebra explicit, proves that empirical averaging and inversion commute, and assigns an
exact supplied-count and product representation to every one of 24 coordinates.  The exact sign
reduction is strong but narrow: it classifies a fixed supplied-count coordinate once its exact
product is known.

The boundary is equally important.  The checked module does not prove component-atom nonnegativity,
does not refine Rust or another executable, does not verify binary64 arithmetic or resource
behavior, and does not make a statistical or higher-source claim.  This document is a bounded
formal audit, not an end-to-end completion certificate for categorical SxPID software or science.

\begingroup
\enlargethispage{0.5\baselineskip}
\footnotesize
\begin{thebibliography}{9}
\setlength{\itemsep}{0.35\baselineskip}
\bibitem{mgw2021}
A.~Makkeh, A.~J. Gutknecht, and M.~Wibral.
Introducing a differentiable measure of pointwise shared information.
\emph{Physical Review E}, 103:032149, 2021.
\url{https://doi.org/10.1103/PhysRevE.103.032149}.

\bibitem{gwm2021}
A.~J. Gutknecht, M.~Wibral, and A.~Makkeh.
Bits and pieces: Understanding information decomposition from part-whole relationships and formal
logic.
\emph{Proceedings of the Royal Society A}, 477, 2021.
\url{https://doi.org/10.1098/rspa.2021.0110}.
\end{thebibliography}
\endgroup

\end{document}
